\documentclass[11pt,a4paper]{article}
\usepackage{fontspec}
\setmainfont{Noto Serif}
\setsansfont{Noto Sans}
\usepackage{babel}
\babelprovide[main,import]{russian}
\usepackage{amsmath,amssymb,mathtools,booktabs,geometry,fancyhdr,microtype,xcolor,hyperref}
\geometry{left=21mm,right=21mm,top=21mm,bottom=21mm}
\definecolor{ink}{HTML}{173D57}
\hypersetup{colorlinks=true,urlcolor=ink,linkcolor=ink,pdftitle={BGS-F. Группа симметрий и квантовые возвраты}}
\pagestyle{fancy}\fancyhf{}\lhead{\small BGS-F}\rhead{\small Группа симметрий и квантовые возвраты}\cfoot{\thepage}
\setlength{\headheight}{15pt}\setlength{\parindent}{1em}\setlength{\parskip}{2.5pt}
\setlength{\emergencystretch}{3em}\tolerance=3000
\newcommand{\Tr}{\operatorname{Tr}}
\newcommand{\EE}{\mathbb E}
\newcommand{\R}{\mathbb R}
\newcommand{\Z}{\mathbb Z}
\newcommand{\C}{\mathbb C}
\newcommand{\si}{\mathrm{Sine}_2}
\newcommand{\eps}{\varepsilon}
\newcommand{\sinc}{\operatorname{sinc}}
\newcommand{\sect}[1]{\section*{\sffamily\color{ink}#1}}
\newcommand{\subt}[1]{\par\medskip\noindent\textbf{#1.}\quad}
\newcommand{\T}{\mathbb T}
\newcommand{\dist}{\operatorname{dist}}
\newcommand{\Ad}{\operatorname{Ad}}
\begin{document}
\noindent{\small\sffamily\color{ink} МАТЕМАТИЧЕСКАЯ ФИЗИКА\hfill 20 НОЯБРЯ 2026}
\begin{center}{\LARGE\sffamily Группа симметрий исходной карты}\\[4pt]
{\large Все итерации и квантовые возвраты к каноническим картам}
\end{center}
Пусть $F=\phi^\tau B$ --- классическая карта исходной задачи, $\tau=10^{-80}$, $\phi$ --- поток её физического символа, а $B=\left(\begin{smallmatrix}3&-4\\-2&3\end{smallmatrix}\right)$. Пусть $\lambda$ --- нормированная площадь тора.

\textbf{Полная классификация.} Для каждого целого $n\ge1$ и каждого сохраняющего площадь гомеоморфизма $\Psi$ тора
\[
       \boxed{\ \Psi F^n\Psi^{-1}=F^{\epsilon n},\quad
             \epsilon\in\{+1,-1\}
             \quad\Longrightarrow\quad
             \epsilon=+1,\quad \Psi=F^j,\ j\in\Z.\ }                \tag{1}
\]
Таким образом, у всех итераций одна и та же группа сохраняющих площадь симметрий, и ни одна итерация не имеет сохраняющего площадь обращения времени. Более того, всякий топологический нормализатор, не входящий в эту группу, переводит площадь во взаимно сингулярную меру. В частности, классификация охватывает все липшицевы нормализаторы и все $C^1$-диффеоморфизмы. Это результат для неизменной силы, без ограничения на $n$.

\textbf{Квантовое следствие.} Пусть $\kappa$ --- фиксированная аналитическая симплектическая карта, а $Q_N$ --- любая её унитарная реализация, удовлетворяющая операторной формуле Егорова на фиксированных гладких наблюдаемых. Если $\kappa$ не является степенью $F$, то для исходных матриц $U_N=e^{i\tau NA_N}M_N$
\[
            \boxed{\quad
       \sup_{m\in\Z}\frac{|\Tr(Q_N^*U_N^m)|}{N}\longrightarrow0 .
             \quad}                                               \tag{2}
\]
В частности, шаблоном может быть любая фиксированная нетождественная аффинная симплектическая карта. В отличие от предшествующей оценки для сдвигов, (2) допускает произвольные фиксированные аналитические нелинейные шаблоны. Равномерность относится к квантовому времени $m$, которое может быть порядка $N$ или больше. При количественной формуле Егорова получена скорость $O_\kappa((\log\log\log N)^{-1/2})$.

Механизм (1) использует две ненулевые суммы Фурье вдоль дуальных гиперболических орбит. Они переводят возможную трансляцию из множества периодических сдвигов $B^n$ в четыре неподвижных сдвига $B$. Затем точные мультипликаторы четырёх неподвижных точек и мнимые части тех же сумм исключают оставшиеся варианты.

\textbf{Граница.} Обычный скалярный возврат соответствует шаблону $\kappa=\mathrm{id}=F^0$ и остаётся открытым при фиксированном кике. Теорема (1) классифицирует классические симметрии. Она сама по себе не устанавливает локальный закон собственных фаз.
\newpage
\sect{1. Топологические нормализаторы линейной карты}
Пусть $H$ --- близкое к тождеству гомеоморфное сопряжение $F H=HB$, существование и количественная регулярность которого ниже проверяются. При $\Psi F^n\Psi^{-1}=F^{\epsilon n}$ карта $A=H^{-1}\Psi H$ удовлетворяет $A B^n=B^{\epsilon n} A$. У неё есть поднятие на плоскость вида $Cx+u(x)$, где $C\in\mathrm{GL}(2,\Z)$, $u$ периодична. На гомологиях получаем $CB^n=B^{\epsilon n}C$. На поднятиях $u(B^n x)-B^{\epsilon n}u(x)$ есть фиксированный целочисленный вектор. Вычитаем его постоянное решение $b$ для $u$ и получаем
\[
                 (u-b)(B^n x)=B^{\epsilon n}(u-b)(x).
\]
Так как $u-b$ ограничена во всех положительных и отрицательных итерациях, гиперболичность вынуждает $u-b=0$. Значит, $A(x)=Cx+b$ на торе, $B^n b=b\bmod\Z^2$.

Действительные собственные значения $B$ равны $\Lambda=3+2\sqrt2$ и $\Lambda^{-1}$. У $B^n$ те же собственные прямые, поэтому $C$ сохраняет их в случае $\epsilon=1$ и меняет местами в случае $\epsilon=-1$. Отсюда уже $CB=B^\epsilon C$. Пусть $R=\operatorname{diag}(1,-1)$, $L=\left(\begin{smallmatrix}1&-2\\-1&1\end{smallmatrix}\right)$. Имеем $L^2=B$ и $RBR=B^{-1}$. В случае $\epsilon=1$ прямое решение $CB=BC$ даёт
\[
                 C=\begin{pmatrix}a&-2b_0\\-b_0&a\end{pmatrix},
              \qquad a^2-2b_0^2=\pm1 .
\]
Все такие целые единицы равны $\pm(1+\sqrt2)^s$, $s\in\Z$. Проверка проста: после знака и умножения на степень $1+\sqrt2$ действительное значение единицы можно поместить в $[1,1+\sqrt2)$. При норме $+1$ разность значения и его сопряжённого меньше $2\sqrt2$, так что целый коэффициент при $\sqrt2$ равен нулю. При норме $-1$ сумма значения и его обратного тоже строго меньше $2\sqrt2$, что несовместимо с положительным целым коэффициентом при $\sqrt2$. Следовательно, единственная приведённая единица --- 1. В случае $\epsilon=-1$ применяем эту классификацию к $CR$.

Итак, без каких-либо условий на меру все топологические нормализаторы имеют вид
\[
           A(x)=\pm L^s R^\eta x+b,\qquad
           s\in\Z,\quad \eta\in\{0,1\},\quad B^n b=b .
                                                                    \tag{3}
\]
В (3) равенство для $b$ понимается на торе, а $\epsilon=(-1)^\eta$. Остаётся выяснить, какие из (3) сохраняют перенесённую площадь.
\newpage
\sect{2. Площадь и уравнение на суммы вдоль орбит}
Пусть $D F(Hx)r_\tau(x)=\exp(f(x))\,r_\tau(Bx)$, где $r_\tau(x)$ задаёт нормированную неустойчивую прямую. Ниже доказывается, что $f$ гёльдерова. Отличие $f$ от логарифма неустойчивого якобиана, перенесённого через $H$, есть кограница: это разность логарифмов норм $r_\tau$ в соседних точках. Для степени используем $f_n(x)=\sum_{a=0}^{n-1}f(B^a x)$.

Положим $\mu=(H^{-1})_*\lambda$. Площадь есть равновесная мера смешивающей карты Аносова $F^n$ для минус логарифма её неустойчивого якобиана. Это стандартное следствие формулы энтропии Песина и неравенства Рюэлля или теоремы о мере Синая--Рюэлля--Боуэна. Сопряжение сохраняет вариационный принцип, а добавление кограницы сохраняет равновесную меру. Значит, $\mu$ есть равновесная мера $B^n$ для $-f_n$.

Если $\Psi$ сохраняет площадь, то $A$ из (3) сохраняет $\mu$. Отображение $A$ биективно переводит множество инвариантных мер $B^n$ в себя. В случае обращения времени используется равенство энтропий $B^n$ и $B^{-n}$ для инвариантной меры. Поэтому давления $-f_n$ и $-f_n\circ A$ одинаковы, а $\mu$ является равновесной мерой обоих потенциалов.

Здесь нужен точный вывод, а не статистическое сравнение. Гиббсовы оценки этой же меры для двух гёльдеровых потенциалов, применённые к одному шару Боуэна малого фиксированного радиуса, ограничивают разность их сумм за любое время единой константой. На орбите периодической точки длины $r$ берём время $r\ell$ и делим на $\ell\to\infty$. Получаем нулевую сумму $f_n\circ A-f_n$ на каждой периодической орбите $B^n$. Теорема Лившица даёт непрерывную функцию $g$ с
\[
                         f_n\circ A-f_n=g\circ B^n-g.              \tag{4}
\]

Пишем $P=B^{\mathsf T}$ и для ненулевой целой метки $\nu$ вводим
\[
                  \mathcal I_\nu=\sum_{q\in\Z}\widehat f(P^q\nu).
\]
Ряд абсолютно сходится: коэффициенты гёльдеровой функции убывают как отрицательная степень частоты, а частоты растут экспоненциально в обе стороны. Ненулевая целая метка не лежит на иррациональной собственной прямой. Суммируя коэффициенты (4) вдоль $P^n$-орбиты $\nu$, правая часть телескопируется в ноль по лемме Римана--Лебега. В левой части сумма по $n$ слагаемым $f_n$ восстанавливает всю $P$-орбиту. Фаза трансляции постоянна на рассматриваемой $P^n$-орбите, поскольку $B^n b=b$. Таким образом, для $\eta=C^{-\mathsf T}\nu$ получаем необходимое тождество
\[
                      \boxed{\ \mathcal I_\nu
                      =e^{2\pi i\eta\cdot b}\,\mathcal I_\eta .\ }   \tag{5}
\]
Оно верно при обеих ориентациях времени. В случае $\epsilon=-1$ порядок прохода по орбите меняет знак, а двухсторонняя сумма остаётся той же.
\newpage
\sect{3. Количественная регулярность сопряжения}
Приведём достаточные границы, чтобы первое слагаемое разложения не было формальным при заданном $\tau$. В координатах $(q,p)$
\[
 g=\cos14\pi q-\tfrac1{10}
     +r(q)\bigl(\tfrac1{10}+\tfrac1{100}\sin2\pi q\bigr)
     +\tfrac1{100}r(q+\tfrac12)\cos(2\pi p-\alpha_0).
\]
Коэффициенты физического ряда получаются из исходных вейлевых коэффициентов умножением на $(-1)^{ab}$ и перестановкой частот $(a,b)\mapsto(b,a)$. Векторное поле $Y=-J\nabla g/(2\pi)$, $J=\left(\begin{smallmatrix}0&1\\-1&0\end{smallmatrix}\right)$. Из носителя $|\nu|_1\le10$ и суммы модулей коэффициентов меньше $6/5$ получаем $\max_{0\le j\le4}\|D^jY\|_\infty<10^{10}$ для операторных норм по максимальной норме на векторах. Достаточно оценки $12(20\pi)^4<10^{10}$. Обозначим эту грубую константу через $K=10^{10}$.

Пусть $d_s=\phi^s-\mathrm{id}$ в поднятии. При $|s|\le10^{-80}$ в дифференцированных уравнениях потока получаем $\|D\phi^s\|\le2$, $\|D^2\phi^s\|\le8K|s|$, $\|D^3\phi^s\|\le100K|s|$. Например, второй порядок имеет источник $D^2Y(D\phi,D\phi)$ и линейный член $DY D^2\phi$, а третий добавляет $D^3Y(D\phi,D\phi,D\phi)$ и три члена $D^2Y(D\phi,D^2\phi)$. Интегрирование с линейным множителем $\le2$ даёт указанные границы. Следовательно,
\[
            \|d_s\|_{C^3}\le10^{12}|s|,
      \qquad \|d_s-sY\|_{C^2}\le10^{26}s^2 .                         \tag{6}
\]
Для второго неравенства интегрируем $Y\circ\phi^v-Y$ по $v$. Внутри пользуемся $Y(x+d_v)-Y(x)=\int_0^1 DY(x+\rho d_v)d_v\,d\rho$. Норма $DY(x+\rho d_v)$ в $C^2$ меньше $10K$ по цепному правилу, а произведение в $C^2$ стоит множитель не более 4. Эти константы оставляют запас в (6).

Возьмём $\beta=1/4$ и норму $C^\beta$ как сумму супремума и семинормы. Для $\mathcal L h=h-B h\circ B^{-1}$ обратный на устойчивой части есть $\sum_{j\ge0}B^j P_s f\circ B^{-j}$, на неустойчивой --- $-\sum_{j\ge1} B^{-j}P_u f\circ B^j$. Проекторы имеют норму $<2$, $7^\beta<2$, $\Lambda>5$, поэтому $\|\mathcal L^{-1}\|_{C^\beta\to C^\beta}<10$. Уравнение $F H=HB$ для $H=\mathrm{id}+h$ равносильно
\[
                    \mathcal L h=d_\tau(\mathrm{id}+B h\circ B^{-1}).
\]
Оператор $h\mapsto Bh\circ B^{-1}$ имеет норму $<14$. На шаре $\|h\|_\beta\le10^{15}\tau$ правая часть до $\mathcal L^{-1}$ имеет норму $<3\cdot10^{12}\tau$ по (6). Для разности подставляемых аргументов применяем интегральную формулу с $Dd_\tau$: семинорма произведения оценивается суммой двух стандартных произведений норм, а нормы обоих периодических добавок меньше 1. Это даёт константу Липшица не больше $10^6\|d_\tau\|_{C^2}$ до $\mathcal L^{-1}$. Получаем сжатие с константой меньше $10^{20}\tau<1$ и решение с $\|h\|_\beta\le10^{15}\tau$.

Карта $H$ сюръективна по степени один. Если $H x=H y$, сопряжение даёт $\operatorname{dist}(B^j x,B^j y)\le2\|h\|_\infty$ для всех целых $j$. Для разностей, центрированных на торе, радиус $1/100$ запрещает решёточные скачки, так как $\|B\|_\infty\le7$. Ограниченность обеих линейных итераций вынуждает разность нуль. Поэтому $H$ инъективна и является гомеоморфизмом.
\newpage
\sect{4. Две ненулевые суммы Фурье при фиксированной силе}
Из (6) и $\|h\|_\beta\le10^{15}\tau$ получаем
\[
       D F(Hx)=[I+\tau DY(Bx)]B+E(x),\qquad
       \|E\|_\beta\le10^{40}\tau^2.                                \tag{7}
\]
Действительно, $DF(Hx)=(I+Dd_\tau(BHx))B$. Ошибка замены $Dd_\tau$ на $\tau DY$ контролируется $C^1$-нормой из (6), а замена $DY(BHx)$ на $DY(Bx)$ стоит не больше $10^4\|Y\|_{C^3}\|h\|_\beta$ по интегральной формуле разности. Константа $10^4$ покрывает обе композиции с $B$ и нормы произведения. По тем же оценкам $\|DF\circ H-B\|_\beta\le10^{15}\tau$.

В базисе $r_u=(1,-1/\sqrt2)$, $r_s=(1,1/\sqrt2)$ матрица кокцикла имеет вид $\left(\begin{smallmatrix}a&b\\c&d\end{smallmatrix}\right)$, где отклонение от $\operatorname{diag}(\Lambda,\Lambda^{-1})$ в $C^\beta$ меньше $10^{16}\tau$. Операция на наклонах $z$ есть $[(c+dz)/(a+bz)]\circ B^{-1}$. В шаре $\|z\|_\beta\le10^{18}\tau$ обратная величина к знаменателю имеет норму меньше $1/5$ по геометрическому ряду в банаховой алгебре $C^\beta$. Числитель по норме меньше $0.22\cdot10^{18}\tau$, а композиция стоит множитель меньше 2. Разность двух дробей выражается через $(ad-bc)(z-z')$ и произведение знаменателей. Так как $ad-bc=1$, константа сжатия не больше $2/25$. Получаем инвариантную неустойчивую прямую со столь малым наклоном.

По определению $f=\log(a+bz)$. Пусть $\ell_u=(1/2,-1/\sqrt2)$ --- левый вектор с $\ell_u r_u=1$. Разложение (7) даёт
\[
       f=\log\Lambda+\tau f_1+\mathcal R,\qquad
       f_1(x)=\ell_u DY(Bx)r_u,\qquad
       \|\mathcal R\|_\beta\le10^{45}\tau^2 .                       \tag{8}
\]
Здесь $bz=O_{C^\beta}(10^{34}\tau^2)$, а квадратичный остаток логарифма ограничивается степенным рядом в той же норме. Явное вычисление производной даёт
\[
 f_1=-\frac{g_{qq}-g_{pp}/2}{2\pi\sqrt2}\circ B .
\]
Обе дуальные орбиты $e_1=(1,0)$ и $e_2=(0,1)$ пересекают носитель $\widehat g$ ровно в исходной точке. При ненулевой степени $P$ модуль второй координаты $P^j e_1$ не меньше 4, а у $P^j e_2$ --- не меньше 3. В исходном ряду $|\nu_2|\le1$. Сдвиг частоты на $P$ из-за композиции с $B$ не меняет орбитальной суммы. Поэтому
\[
  \mathcal I_{e_1}=\tau\sqrt2\pi\,\frac{200-20i}{4000}
                      +O(20\cdot10^{45}\tau^2),\qquad
  \mathcal I_{e_2}=-\frac{\tau\sqrt2\pi}{2}\,
                    \frac{-12-16i}{4000}
                      +O(20\cdot10^{45}\tau^2).                   \tag{9}
\]
Для остатка используем $|\widehat{\mathcal R}(\nu)|\le\|\mathcal R\|_\beta|\nu|_\infty^{-\beta}$, полученное сдвигом на половину периода максимальной координаты. Диагональные элементы $P^j$ дают $|P^j e_i|_\infty\ge\Lambda^{|j|}/2$. Поэтому сумма этих весов меньше 20. При $\tau=10^{-80}$ обе мнимые части в (9) по модулю превышают $10^{-3}\tau$.
\newpage
\sect{5. Устранение всех периодических сдвигов}
В (5) заменим $\nu$ на $P\nu$. Тогда $\eta$ заменяется на $P^\epsilon\eta$, поскольку $CB=B^\epsilon C$. Ни одна из двух орбитальных сумм от этого не меняется. Сравнение тождеств даёт при $\mathcal I_\eta\ne0$
\[
                         \eta\cdot(B^\epsilon-I)b\in\Z .
\]
Берём $\eta=e_1,e_2$, что допустимо по (9) и обратимости $C$. Получаем $Bb=b$ на торе, хотя в (3) было известно только $B^n b=b$. Следовательно, $A B=B^\epsilon A$, то есть исходная $\Psi$ уже нормализует саму карту $F$.

Четыре неподвижные точки $B$ суть $(q,p)=(a/2,b_0/2)$, $a,b_0\in\{0,1\}$: определитель $B-I$ по модулю равен 4, а все четыре указанные точки неподвижны. Все они также неподвижны у $F$. В них $\nabla g=0$, а $g_{qp}=g_{pp}=0$, поскольку множитель $r(q+1/2)$ при $p$-зависимом члене имеет двойной нуль. Вычисление оставшейся производной даёт в порядке $00,01,10,11$
\[
       \frac{g_{qq}}{\pi^2}=\frac1{125}
              (-19606,\ -19594,\ 24306,\ 24894).
\]
Например, каждый элемент получается из $-4\sum_{j,l}j^2\widehat g(j,l)(-1)^{ja+lb_0}$. Матрица $DY$ в этих точках имеет единственный возможный ненулевой элемент ниже диагонали. Поток в точке неподвижен, а $D\phi^\tau=I+\tau DY$ точно. Поэтому
\[
                    \Tr DF = 6-\frac{2\tau}{\pi}g_{qq}.
\]
Все четыре следа различны и больше 2. При определителе один это даёт четыре разных неустойчивых мультипликатора. Сопряжение $H$ сопоставляет точки с теми же координатами благодаря малой равномерной норме $h$.

У неподвижной точки уравнение (4) даёт $f_n(Ax)=f_n(x)$, то есть $f(Ax)=f(x)$, так как после $Bb=b$ обе точки неподвижны у $B$. Значение $f$ в них точно логарифм неустойчивого мультипликатора. Поэтому $A$ фиксирует каждую из четырёх точек. Значит, $b=0$ и $C=I\bmod2$. В (3) это означает чётность $s$, так как $L\bmod2$ имеет порядок 2, а $R$ и знак равны $I\bmod2$. Остались $C=\pm B^j R^\eta$.

Если $C=-B^j$, то (5) при $b=0$ и $\nu=e_1$ даёт $\mathcal I_{e_1}=\mathcal I_{-e_1}=\overline{\mathcal I_{e_1}}$, что противоречит (9). Если $C=B^jR$, тот же аргумент действует для $e_2$. Если $C=-B^jR$, действует для $e_1$. Факторы $B^j$ в этих сравнениях меняют только представителя той же орбиты. Следовательно, остаются $C=B^j$, $\epsilon=+1$. Возвращаясь через $H$, получаем $\Psi=F^j$. Все такие степени действительно сохраняют площадь и коммутируют. Это доказывает (1) одновременно для каждого $n\ge1$.
\newpage
\sect{6. Сингулярность остальных нормализаторов}
Для каждой итерации $F^n$ площадь эргодична. Например, по сопряжению с топологически смешивающим $B^n$ она переносится в равновесную меру гёльдерова потенциала, которая эргодична по теореме Боуэна. Если $\Psi F^n\Psi^{-1}=F^{\epsilon n}$, то $\Psi_*\lambda$ также инвариантна и эргодична для $F^n$, при любом знаке $\epsilon$. Различные эргодические меры для одного преобразования взаимно сингулярны. Это можно увидеть по непересекающимся множествам типичных точек для средней по времени одной непрерывной функции, интегралы которой в этих мерах различны.

Если $\Psi_*\lambda=\lambda$, то уже применимо (1). Во всех остальных случаях $\Psi_*\lambda\perp\lambda$. Это, в частности, исключает для остальных нормализаторов свойство Лузина $N$ в любом направлении. Действительно, при сингулярности есть множество $E$ полной площади, для которого $\Psi(E)$ имеет площадь ноль. Если $\Psi$ переводит нуль-множества в нуль-множества, то и $\Psi(E^c)$ имеет площадь ноль, что противоречит сюръективности. Если это свойство имеет $\Psi^{-1}$, то из нулевой площади $\Psi(E)$ следовала бы нулевая площадь $E$. Липшицевы отображения имеют свойство $N$, как и $C^1$-диффеоморфизмы на компактном торе. Поэтому (1) охватывает в частности все такие нормализаторы. Отсутствие сохраняющего площадь обращения не зависит от гладкости предполагаемой симметрии.

\subt{Следствие для фиксированного квантового времени} Пусть $Q_N$ унитарна или антиунитарна и удовлетворяет формуле Егорова для фиксированной гладкой канонической карты $\kappa$ на вещественных наблюдаемых. Унитарному выбору соответствует симплектическая $\kappa$, антиунитарному --- антисимплектическая. Для фиксированного $n\ge1$ и знака $\epsilon$ положим $X_N=Q_N U_N^nQ_N^{-1}$, $Y_N=U_N^{\epsilon n}$. Если $\kappa F^n\kappa^{-1}\ne F^{\epsilon n}$, то
\[
        \liminf_{N\to\infty}\inf_{|z|=1}
                              \|X_N-zY_N\|_{\rm n}>0,              \tag{10}
\]
где $\|A\|_{\rm n}^2=N^{-1}\Tr A^*A$. Действительно, найдём вещественный гладкий символ $a$, перенос которого двумя классическими картами различается. Формула Егорова для $X_N,Y_N$ приводит разность сопряжений наблюдаемого $\operatorname{Op}_N(a)$ к квантованию этой ненулевой разности. Нормированная норма Гильберта--Шмидта такого квантования стремится к его норме $L^2$ по ортогональности операторов Вейля и убыванию коэффициентов Фурье. С другой стороны, разность сопряжений имеет норму $\le 2\|\operatorname{Op}_N(a)\|_{\rm op}\|X_N-zY_N\|_{\rm n}$. Верхняя граница для $\|\operatorname{Op}_N(a)\|_{\rm op}$ дана суммой модулей Фурье. Это доказывает (10). В частности, для антиунитарного выбора (1) запрещает совпадение классических карт при обоих знаках: никакая степень $F$ не антисимплектична. Здесь фиксируется $n$, как и в формуле Егорова. Следующая оценка по квантовому времени имеет другой вид и равномерна по всем $m$.
\newpage
\sect{7. След при несовпадающем аналитическом шаблоне}
Установим общую лемму. Пусть $V_N$ --- последовательность унитарных операторов с формулой Егорова для аналитической симплектической карты $G$ на торе. Если $G\ne\mathrm{id}$, то $\Tr V_N/N\to0$. Возьмём конечное среднее
\[
       V_{N,L}=L^{-2}\sum_{0\le k_1,k_2<L}
                                 T_N(k)^*V_N T_N(k).
\]
След среднего тот же. По Коши $|\Tr V_N/N|^2\le\|V_{N,L}\|_{\rm n}^2$. Раскрывая квадрат и используя вейлево произведение, получаем взвешенное среднее нормированных следов $V_N^*T_N(k-k')V_N T_N(k-k')^*$. Фазы, относящиеся лишь к выбору операторов Вейля, сокращаются. При фиксированном $L$ формула Егорова переводит эти конечные следы в интегралы соответствующих характеров от $G(x)-x$. Предельная норма квадрата имеет формулу
\[
 \int_{\T^2}
      \prod_{\ell=1}^2
        \left|L^{-1}\sum_{j=0}^{L-1}
                 e^{2\pi i j(G(x)-x)_\ell}\right|^2\,dx .
\]
Если в Егорова используется обратная карта, замена переменной даёт тот же интеграл. Перестановка координат и смена знаков в симплектическом спаривании также не меняют произведение. Это та же конечная процедура усреднения, которая использовалась для сдвигов. При $L\to\infty$ произведение стремится к характеристической функции множества неподвижных точек. У аналитической нетождественной $G$ это множество меры нуль: хотя бы одна глобальная аналитическая функция $e^{2\pi i(G(x)-x)_\ell}-1$ не равна нулю тождественно. По доминированной сходимости интеграл стремится к нулю. Два предела берутся именно в этом порядке.

Теперь пусть $Q_N$ имеет аналитическую каноническую карту $\kappa$ из формулировки (2). Для $j=0,\ldots,L-1$ положим $Q_{N,j}=U_N^j Q_N U_N^{-j}$. Их классические карты суть $\kappa_j=F^j\kappa F^{-j}$. При $i\ne j$ совпадение $\kappa_i=\kappa_j$ означало бы, что $\kappa$ коммутирует с ненулевой степенью $F$. По (1) это возможно только для степени $F$, а она исключена условием теоремы. Поэтому все $\kappa_j$ различны. Лемма для $Q_{N,i}^*Q_{N,j}$ даёт предельную ортогональность в нормированном гильбертовом скалярном произведении. Следовательно,
\[
           \left\|\frac1L\sum_{j=0}^{L-1}Q_{N,j}\right\|_{\rm n}^2
                                                    =L^{-1}+o(1).
\]
Для каждой степени $U_N^m$ все скалярные произведения с $Q_{N,j}$ в точности совпадают, так как степень коммутирует с $U_N$. По Коши
\[
       \limsup_{N\to\infty}\sup_{m\in\Z}
            |\Tr(Q_N^*U_N^m)|/N\le L^{-1/2}.
\]
Увеличивая $L$, получаем (2). Здесь формула Егорова нужна только для фиксированных $j$ и фиксированного числа композиций. Само время $m$ в Егорова не подставляется.
\newpage
\sect{8. Применения и граница результата}
Ни одна ненулевая степень $F^j$ не является аффинной. Действительно, всякое аффинное отображение $F^j$ переводило бы непрерывную группу сдвигов в себя при сопряжении. В части о жёсткости Вейля доказано для этого же фиксированного $\tau$, что $F^j t_a F^{-j}=t_b$ возможно лишь при $a=b=0$. Отрицательные $j$ сводятся к положительным обращением равенства. Поэтому любая фиксированная нетождественная аффинная каноническая карта подходит в (2). В частности, для каждой фиксированной ненулевой степени чистого линейного кванта $M_N^r$ имеем
\[
           \sup_{m\in\Z}N^{-1}|\Tr((M_N^r)^*U_N^m)|\longrightarrow0.
\]
Это распространяется на аналитические нелинейные шаблоны и не требует простой спектральности. Для унитарных шаблонов из (2) также следует, что минимальное по скалярной фазе расстояние $\|U_N^m-zQ_N\|_{\rm n}$ равномерно по $m$ стремится к $\sqrt2$.

Если некоторая последовательность степеней $U_N^{m_N}$ сходится по действию Егорова на фиксированных наблюдаемых к гладкой канонической карте $\kappa$, то $\kappa$ коммутирует с $F$ по точному коммутированию матриц. Следовательно, $\kappa=F^r$ по (1). Это ограничение относится к существованию канонического предела, а не устанавливает существование такого предела при произвольном выборе времени.

Классификация и квантовые оценки покрывают прежние вопросы о возможных нетривиальных гладких симметриях. Они оставляют открытыми скалярные пики $\Tr U_N^m$ на бесконечных последовательностях гейзенберговских времён: у шаблона $\kappa=\mathrm{id}$ все члены среднего совпадают. Для полной гипотезы BGS-F не вычислены обычный форм-фактор и локальные корреляции при исходном фиксированном кике.

\subt{Воспроизведение} Код проверяет все использованные гауссовы рациональные коэффициенты из точного символа, четыре следа производных неподвижных точек, алгебру $L,R,B$, остаток в (9) и отсечение вариантов по симметриям. Проверка не заменяет аналитические оценки потока, сопряжения или термодинамического формализма. Они приведены полностью с численным запасом при $10^{-80}$.

\subt{Стандартные результаты} Для равновесных мер и геометрического потенциала используются теоремы Боуэна и Синая--Рюэлля--Боуэна в форме, изложенной, например, во введении к работе В. Клименхаги, Т. Фишера и Д. Томпсона \href{https://www.math.uh.edu/~climenha/doc/CFT-Mane.pdf}{Equilibrium states for Mañé diffeomorphisms}. Нужная аддитивная теорема Лившица приведена в A. Wilkinson, \href{https://www.math.uchicago.edu/~wilkinso/papers/livsic.pdf}{The cohomological equation for partially hyperbolic diffeomorphisms}, теорема 0.1. Для количественных границ здесь применяются элементарные ряды в фиксированной гёльдеровой норме и точная арифметика исходного символа.
\newpage
\sect{9. Растущее окно для аналитических шаблонов}
Для шаблонов с количественной формулой Егорова можно усилить (2) до скорости. Пусть для $Q_N$ и его обратного ошибка Егорова в операторной норме не больше $C_\kappa N^{-1}\|a\|_{C^7}$ для гладких символов $a$. Это в частности выполняется для фиксированного числа линейных метаплектических шагов и гладких гамильтоновых киков. Тогда для $\kappa$ из (2) получаем
\[
       \sup_{m\in\Z} N^{-1}|\Tr(Q_N^* U_N^m)|
                  =O_\kappa\bigl((\log\log\log N)^{-1/2}\bigr).     \tag{11}
\]

Ключевой факт состоит в равномерной отделённости $\kappa$ от $F^j\kappa F^{-j}$ по всем $j\ge1$. Для доказательства пишем $H^{-1}\kappa H=Cx+u(x)$ на поднятии. Если $C$ не коммутирует с $B$, то $B^j C B^{-j}\ne C$ для всякого положительного $j$. Получившиеся карты имеют разные гомологии и не могут быть ближе инъективного радиуса равномерной метрики тора. Если $C$ коммутирует с $B$, разность периодических частей двух поднятий равна $B^j u(B^{-j}x)-u(x)$. При непостоянной неустойчивой компоненте $u$ её осцилляция здесь не меньше $(\Lambda^j-1)$ исходной. Если меняется только устойчивая компонента, осцилляция не меньше $(1-\Lambda^{-j})$ исходной. Малая равномерная разность на торе ограничивала бы осцилляцию поднятой разности, поскольку выбор ближайшего решёточного вектора был бы постоянным. Обе оценки дают отделённость, если $u$ не постоянна.

Если $u=b$ постоянна, аффинная карта $Cx+b$ сохраняет $\mu=(H^{-1})_*\lambda$ и коммутирует с $B$ на гомологиях. Её коммутатор с $B$ есть сохраняющий $\mu$ сдвиг на $(B-I)b$. Группа всех сохраняющих $\mu$ сдвигов замкнута и инвариантна относительно $B$. Она тривиальна: положительноразмерная связная компонента дала бы рациональное инвариантное направление $B$ или весь тор, а всякий элемент конечной группы периодичен под $B$. В случае всего тора берём ненулевой полурешёточный сдвиг. В обоих последних случаях ненулевой периодический сдвиг противоречит (1), будучи нормализатором некоторого $B^n$ через $H$. Поэтому $Bb=b$ и исходная $\kappa$ была бы степенью $F$, что исключено. Перенос нижней границы через $H$ использует лишь равномерную непрерывность $H^{-1}$.

Значит, для $G_j=\kappa^{-1}F^j\kappa F^{-j}$ величина $\sup_x\operatorname{dist}(G_jx,x)$ положительно ограничена снизу по всем $j\ge1$. Карты аналитичны в комплексной трубке радиуса $cK^{-j}$ с равномерной границей на мнимую часть. Это следует итерацией оценки мнимой части каждого из фиксированных множителей через его первую производную на меньшей трубке. Применяем лемму об аналитической полосе из части о возвратах со сдвигом, где в разделе 2 она доказана формулой Йенсена и факторизацией Бляшке. Получаем
\[
 \operatorname{area}\{x:\operatorname{dist}(G_jx,x)<\epsilon\}
       \le C e^{CK^j}\epsilon^{\exp(-CK^j)}.
\]
Вейлево усреднение из раздела 7 можно теперь вести до $L=\lfloor N^{1/24}\rfloor$. Ошибка составной формулы Егорова не больше $C^{j+1}N^{-1}(1+L)^{12}$: норма переноса в $C^7$ умножается на фиксированную константу за шаг, а ошибка произведения и следа контролируется вейлевым правилом, как в разделе 3 части о сдвигах. Вне области $\operatorname{dist}(G_jx,x)<L^{-1/2}$ геометрическая сумма из раздела 7 даёт $O(1/L)$. Выбирая $T=\lfloor c_\kappa\log\log\log N\rfloor$ с достаточно малым $c_\kappa$, получаем $|\Tr(Q_N^*U_N^j Q_N U_N^{-j})|/N\le e^{-(\log N)^{3/4}}$ при $1\le j\le T$. Среднее нормализованных операторов $U_N^j Q_NU_N^{-j}$ до $T$ имеет квадрат нормы $O(1/T)$, что доказывает (11) тем же рассуждением по всем $m$.
\newpage
\sect{10. Гладкие веса и спектральные разности}
Квантовая теорема допускает гладкий вес. Пусть $a\in C^\infty(\T^2)$, $D_N=\operatorname{Op}_N(a)Q_N$. При тех же условиях, что в (2), если $\kappa$ не степень $F$, то
\[
       \sup_{m\in\Z}\frac{|\Tr(D_N^*U_N^m)|}{N}\longrightarrow0.     \tag{12}
\]
Более того, пусть $e_i$ --- произвольный ортонормированный собственный базис $U_N$, $\lambda_i$ --- собственные числа. При $a\ne0$ вероятностная мера
\[
       \frac1{N\|D_N\|_{\rm n}^2}
          \sum_{i,j}|\langle e_i,D_Ne_j\rangle|^2
                                \delta_{\lambda_i\overline{\lambda_j}}
\]
слабо сходится на окружности к равномерной мере. В отличие от обычного форм-фактора, здесь спектральные разности взвешены переходами конкретного оператора.

Для доказательства сначала усилим лемму раздела 7: если $V_N$ реализует аналитическую нетождественную карту $G$, то $N^{-1}\Tr(\operatorname{Op}_N(b)V_N)\to0$ для любого фиксированного гладкого $b$. Действительно, для каждого фиксированного вейлева оператора $T_N(v)$ его классический сдвиг имеет размер $O(1/N)$, и $T_N(v)V_N$ реализует ту же карту $G$. При конечной сумме Фурье применяем лемму к каждому слагаемому. За конечной суммой ошибка ограничена хвостом абсолютно сходящегося ряда. Это даёт заявленное усиление.

Сопряжённые операторы $D_{N,j}=U_N^jD_NU_N^{-j}$ имеют вид гладкого веса перед $Q_{N,j}$ с ошибкой $o(1)$ при фиксированном $j$. Скалярное произведение $D_{N,i},D_{N,j}$ получается перестановкой множителей из нормированного следа $Q_{N,i}^*Q_{N,j}$ со вновь гладким весом. При $i\ne j$ применима усиленная лемма, как в доказательстве (2). При $i=j$ норма квадрата стремится к $\|a\|_{L^2}^2$. Поэтому среднее по $L$ сопряжениям имеет предельный квадрат нормы $\|a\|_{L^2}^2/L$. Коммутирование $U_N^m$ с $U_N$ снова даёт (12) при $L\to\infty$. Моменты выписанной меры являются этими же гильбертовыми корреляциями при $i=0$. Все ненулевые фиксированные моменты исчезают, что даёт равномерный предел.

Если же $Q_N=U_N^r$ при фиксированном целом $r$, то
\[
 \sup_{m\in\Z}\frac1N\left|
    \Tr\big((\operatorname{Op}_N(a)U_N^r)^*U_N^m\big)
           -\overline{\int a\,d\lambda}\,\Tr U_N^{m-r}\right|\to0.  \tag{13}
\]
Вычитаем среднее из $a$. Усредняем его квантование под сопряжением до фиксированного $L$. Предельная нормированная норма равна $L^2$-норме классического среднего по $F$, которое исчезает при $L\to\infty$ по эргодичности. Цикличность следа и Коши дают равномерность по $m$. Таким образом, у гладко взвешенных аналитических проб все возможные остаточные макроскопические корреляции сведены к обычным следам со сдвигом времени. Саму величину этого скалярного остатка (13) не вычисляет.
\end{document}
